\subsection{Pilot tiếng Việt dịch và nhãn AI cho yêu cầu hỗ trợ}\label{sec:vi-silver-pilot}
Phép bổ sung này thực nghiệm nhận diện lời yêu cầu hoặc nhu cầu hỗ trợ thiên tai trên \textbf{văn bản dịch sang tiếng Việt và nhãn do AI gán}. Ảnh và văn bản gốc lấy từ cùng bài đăng CrisisMMD~\citep{alam2018crisismmd}; \VietnameseDataReviewVI{} Nhãn request bao gồm lời kêu gọi cứu, sơ tán, cung cấp vật phẩm, quyên góp hoặc nhu cầu hỗ trợ chưa được đáp ứng; không phải mức nguy cấp hay ưu tiên điều phối. Nhãn informative của pilot trước không được chuyển thành request.

\paragraph{Dữ liệu và cách tạo nhãn.} Từ cohort CrisisMMD đã kiểm, chọn 300 cặp theo hash trong từng split sự kiện, một cặp cho mỗi nhóm bản trùng. Một nửa mỗi split được chọn từ các câu có từ khóa rộng về hỗ trợ, nửa còn lại từ các câu khác. Đây là lấy mẫu làm giàu chủ đề, không phản ánh prevalence tự nhiên. AI dịch và gán một nhãn cho từng câu nguồn trong một lượt, kèm lý do; không xem ảnh để gán nhãn. Các lượt được chia thành ba phần để xử lý, không phải ba người gán cùng mẫu và không tạo chỉ số đồng thuận. Bốn câu chưa rõ nghĩa được giữ trong inventory và loại trước fit, không lấy bù. Còn \VISilverTrainN{} train, \VISilverDevN{} validation và \VISilverTestN{} test; test có \VISilverPositiveN{} request và \VISilverNegativeN{} no\_request.

Giữ Sri Lanka làm test, Maria làm validation, năm sự kiện khác làm train; bản sao, ID, hash ảnh và câu dịch được kiểm không giao nhau giữa split. Tập sự kiện này từng được xem trong pilot informative nên đây là bằng chứng phát triển hồi cứu. Bản dịch và nhãn đã được nhóm kiểm tra; chưa có chỉ số đồng thuận gán nhãn độc lập. Vì nhãn dựa vào câu nguồn, phép so sánh đo khớp nhãn văn bản AI của báo cáo, không chứng minh hiểu đúng cả ảnh và câu.

\paragraph{Mô hình đã chạy.} TF-IDF thử word unigram/bigram và char 3--5; vocabulary chỉ fit train. Nhánh văn bản frozen dùng DistilBERT đa ngôn ngữ~\citep{sanh2019distilbert,distilbertMultilingualCard}, mean pooling có attention mask và 768 chiều. Nhánh ảnh dùng MobileNetV3-Large ImageNet frozen, 960 chiều, khác encoder Small của pilot informative và khác classifier mức ngập trong ứng dụng. Các nhánh đơn và fusion frozen chuẩn hóa L2; fusion nối hai block rồi chia $\sqrt{2}$, học LR cân bằng lớp với $C\in\{0{,}1;1;10\}$. Có 18 fit head được chọn bằng macro-F1 validation; ngưỡng giữ 0,5. DistilBERT frozen được chọn làm đối chứng chính từ validation trước test, không chọn lại theo kết quả test.

Nhánh riêng fine-tuning DistilBERT cập nhật encoder và đầu hai lớp trên cùng nhãn AI. Protocol bổ sung khóa trước dự đoán test: hai learning rate $10^{-5}$ và $2\times10^{-5}$, ba seed 17/29/43, tối đa năm epoch, patience 2, batch hiệu dụng 16 và dropout 0,2. Chọn learning rate theo trung bình validation của ba seed, giữ cả ba checkpoint tương ứng; không chỉ báo seed có test tốt nhất. Bảng kết quả dùng lần tái chạy CUDA ngày 10/11/2026 trên RTX 5060 Laptop, chọn $2\times10^{-5}$; kiểm tra một tensor backbone ở mỗi checkpoint xác nhận trọng số thay đổi. Lần thực thi trước trên cùng dữ liệu đạt macro-F1 trung bình \VISilverPreviousFineFOne\% và được giữ riêng trong hồ sơ tái lập, không coi hai lần là hai bộ test độc lập. Nhánh này không fine-tuning chung encoder ảnh--văn bản và không thay đổi tương phản chính của fusion frozen.

Bảng~\ref{tab:vi-silver-pilot} trình bày các kết quả đã chạy trên cùng \VISilverTestN{} ID; recall $(+)$ là recall request. Hàng fine-tuning là trung bình ba seed, không tăng số cặp test độc lập lên ba lần. Hai hàng cuối chỉ che một block của fusion cố định, khác chuyển sang classifier đơn.
% Generated by export_vi_silver.py; do not edit.
\begin{table}[htbp]\centering
\caption{Pilot trên văn bản dịch tiếng Việt và nhãn AI: cùng 59 cặp test Sri Lanka, chỉ 6 nhãn request}\label{tab:vi-silver-pilot}
\begin{tabular}{P{.45\textwidth}rrr}\toprule
Cấu hình & Macro-F1 (\%) & Acc. (\%) & Recall (+) (\%) \\ \midrule
Lớp đa số & 47,32 & 89,83 & 0,00 \\
TF-IDF + LR & 45,87 & 84,75 & 0,00 \\
DistilBERT frozen + LR & 44,13 & 66,10 & 16,67 \\
DistilBERT fine-tuning (TB 3 seed) & 50,73 & 78,53 & 16,67 \\
Ảnh frozen + LR & 54,89 & 84,75 & 16,67 \\
TF-IDF + ảnh & 46,36 & 86,44 & 0,00 \\
DistilBERT frozen + ảnh & 64,91 & 88,14 & 33,33 \\
Trung bình xác suất & 56,30 & 86,44 & 16,67 \\
Fusion che ảnh & 47,32 & 89,83 & 0,00 \\
Fusion che văn bản & 59,48 & 83,05 & 33,33 \\
\bottomrule\end{tabular}
\end{table}


Fusion frozen đạt macro-F1 \VISilverFusionFOne\% và recall request \VISilverFusionRecall\%; nó dự đoán đúng hai trong sáu nhãn request, bỏ sót bốn. Lớp đa số đạt Accuracy 89,83\% nhưng recall request bằng không; vì vậy Accuracy riêng lẻ dễ che lỗi bỏ sót trong test lệch lớp. Fine-tuning DistilBERT đạt macro-F1 trung bình \VISilverFineFOne\%, khoảng min--max giữa seed $[\VISilverFineFOneLow;\VISilverFineFOneHigh]$, recall request trung bình \VISilverFineRecall\%; mỗi seed nhận đúng một hoặc hai trong sáu request. Đây là biến động huấn luyện trên cùng dữ liệu, không phải CI qua nhiều quần thể. Điểm T2 thấp hơn fusion frozen trong lần đánh giá này chưa phải một kiểm định chênh lệch giữa hai phương pháp.

Bảng~\ref{tab:vi-silver-effects} dùng 2.000 bootstrap draw ghép cặp theo 59 nhóm test, cùng draw cho mọi cấu hình. Các CI 95\% chưa hiệu chỉnh, có điều kiện trong một sự kiện và đối với nhãn AI; chúng không bao gồm sai số dịch hoặc sai nhãn. Hiệu fusion so với DistilBERT frozen là \VISilverFusionTextDelta{} điểm phần trăm, CI $[\VISilverFusionTextLow;\VISilverFusionTextHigh]$. Các CI đều chứa hoặc chạm 0, chưa xác lập ưu thế fusion. Không áp dụng tương phản này để so trực tiếp lợi ích fine-tuning với fusion có encoder cố định.
% Generated by export_vi_silver.py; do not edit.
\begin{table}[htbp]\centering
\caption{Hiệu macro-F1 fusion frozen so với đối chứng trên nhãn AI; CI 95\% ghép cặp chưa hiệu chỉnh}\label{tab:vi-silver-effects}
\begin{tabular}{P{.42\textwidth}rrr}\toprule
Đối chứng & Hiệu (đpt) & Cận dưới & Cận trên \\ \midrule
DistilBERT frozen (chính) & 20,78 & -3,93 & 42,65 \\
Ảnh frozen & 10,03 & 0,00 & 26,80 \\
Trung bình xác suất & 8,61 & 0,00 & 25,86 \\
\bottomrule\end{tabular}
\end{table}


Khi che ảnh, fusion dự đoán lớp âm cho toàn bộ test và bỏ sót cả sáu request; classifier text đơn dự đoán đúng một request. Đây là dấu hiệu chính sách thiếu nguồn phải được đánh giá riêng, không được gọi việc che block là fallback đã hoạt động. Kết quả được kiểm lại từ dự đoán full precision và công thức confusion/F1 độc lập với hàm tính metric chính. Kiểm tính toán không xác nhận nhãn AI đúng nghiệp vụ. Pilot bổ sung bằng chứng đã chạy NLP, fine-tuning và fusion trên tiếng Việt dịch; đánh giá trên tin thu nhận trực tiếp, nhãn độc lập, trích xuất span và hiệu quả cứu hộ thực địa là bước tiếp theo. Tình trạng tích hợp do nhóm xác nhận được ghi ở Mục~\ref{sec:mobile-integration-status}.
